\chapter{Detaylı Modül İncelemeleri}

% ==========================================
% BÖLÜM 4.1: THISWORKBOOK (ÇALIŞMA KİTABI ETKİLEŞİM KATMANI)
% ==========================================

\section{\texttt{ThisWorkbook}}

\texttt{ThisWorkbook}, Excel VBA mimarisinde çalışma kitabının yaşam döngüsünü (\textit{Lifecycle}) ve global olayları (\textit{Workbook-Level Events}) dinleyen yerleşik sınıf modülüdür. Projede uygulamanın otomatik ilklendirilmesini (\textit{Initialization}) ve bellek değişkenlerinin hazırlanmasını sağlar.

Excel açıldığında grafik motoru, OLE bileşenleri ve Active Document yapısı henüz yüklenmemiş olabilir. Bu süreçte doğrudan makro çalıştırmak veya arayüze OLE nesneleri çizmek; grafik kaymalarına ve Korumalı Görünüm (\textit{Protected View}) modunda uygulamanın aniden çökmesine neden olur.

\texttt{ThisWorkbook} modülü, söz konusu kararsızlığı önlemek adına **Asenkron Gecikmeli Başlatma (Asynchronous Delayed Execution)** mimari desenini uygular:

\begin{itemize}\setlength{\itemsep}{0pt}\setlength{\parskip}{1pt}
    \item \textbf{Yerleşik Olay Yakalama (\texttt{Workbook\_Open}):} Otomatik tetiklenen giriş noktasıdır (\textit{Entry Point}).
    \item \textbf{İşlem Oluğu Sıfırlama (\texttt{DoEvents}):} Bekleyen tüm grafik yenileme ve ekran çizim olaylarını zorunlu kılar.
    \item \textbf{Zamanlanmış Tetikleme (\texttt{Application.OnTime}):} Arayüz kurma fonksiyonu olan \texttt{InitializeAppUI} prosedürünü 1 saniyelik mikro gecikmeyle Excel yürütme kuyruğuna kaydeder.
\end{itemize}

\vspace{0.15cm}

\begin{infobox}{\faClock\ Asenkron Başlatmanın Sistem Performansına Etkisi}
\texttt{Application.OnTime Now + TimeValue("00:00:01"), "InitializeAppUI"} kullanımı, dosya açılışında Excel UI motorunun tamamen stabil hale gelmesini bekler. Bu sayede servisler ve dinamik butonlar çökme riski olmadan sıfır hatayla ilklendirilir.
\end{infobox}

\vspace{0.15cm}

Aşağıdaki kod bloğu, projedeki \texttt{ThisWorkbook} sınıf modülünün tam yapısını içermektedir:

\begin{lstlisting}[caption={ThisWorkbook Sınıf Modülü Kod Yapısı}]
'*****************************************************************************
' Class: ThisWorkbook
' Purpose: Standard class module that handles workbook-level events.
'*****************************************************************************

Private Sub Workbook_Open()
    DoEvents
    Application.OnTime Now + TimeValue("00:00:01"), "InitializeAppUI"
End Sub
\end{lstlisting}

\newpage

% ==========================================
% BÖLÜM 4.2: CLSDYNAMICBUTTON (DİNAMİK NESNE TABANLI OLAY KONTROLÜ)
% ==========================================

\section{\texttt{clsDynamicButton}}

\texttt{clsDynamicButton}, çalışma zamanında (\textit{Runtime}) üretilen butonların tıklanma olaylarını (\texttt{Click Event}) Nesne Tabanlı Programlama (OOP) prensipleri doğrultusunda dinamik olarak yakalayan özel bir Sınıf Modülüdür (\textit{Class Module}).

UserForm kullanılmayan mimarilerde arayüz elemanları doğrudan \texttt{Worksheet} üzerine çizilir. Ancak dinamik eklenen bir \texttt{MSForms.CommandButton} nesnesinin statik bir tıklama prosedürü bulunmaz. Bu kısıtlamayı aşmak adına projedeki her bir dinamik buton, \texttt{clsDynamicButton} sınıfından türetilen bir nesneye bağlanır ve bellek ömrünü uzatmak için global bir koleksiyonda saklanır.

Sınıf içerisindeki \texttt{WithEvents} bildirimi, buton etkileşimlerini anlık dinler. Tıklama anında izlenen işlem adımları:

\begin{itemize}\setlength{\itemsep}{0pt}\setlength{\parskip}{1pt}
    \item \textbf{Referans ve Kapsayıcı Tespiti:} Tıklanan buton örneği yakalanır ve üst paneli (\texttt{MSForms.Frame}) doğrulanır.
    \item \textbf{Olay Delegasyonu:} Panelin adına göre (\texttt{pan\_Login}, \texttt{pan\_OSM} vb.) tetiklenecek yönlendirme \texttt{mod\_UI\_Events} modülündeki handler fonksiyonlarına aktarılır.
\end{itemize}

\vspace{0.15cm}

\begin{warningbox}{\faExclamationTriangle\ Bellek Yönetimi ve Nesne Yaşam Ömrü (Garbage Collection)}
VBA bellek yöneticisi, referansı tutulmayan sınıf örneklerini anında siler. Dinamik butonların tıklama olaylarının aktif kalması için \texttt{clsDynamicButton} örneklerinin \texttt{ButtonCollection} koleksiyonuna \texttt{Add} metoduyla eklenmesi kritik önem taşır.
\end{warningbox}

\vspace{0.15cm}

Aşağıdaki kod bloğu, dinamik buton olaylarını dinleyen ve ilgili UI panellerine yönlendiren sınıf yapısını özetlemektedir:

\begin{lstlisting}[caption={clsDynamicButton Sınıf Modülü Kod Yapısı}]
'*****************************************************************************
' Class: clsDynamicButton
' Purpose: This class is used to handle events for dynamically created buttons.
'*****************************************************************************

Public WithEvents DynamicButton As MSForms.CommandButton

Private Sub DynamicButton_Click()
    Dim parentFrame As MSForms.Frame
    
    On Error Resume Next
    Set parentFrame = DynamicButton.Parent
    On Error GoTo 0

    If parentFrame Is Nothing Then Exit Sub

    Select Case parentFrame.Name
        Case "pan_Login":    mod_UI_Events.HandleLoginPanelEvents DynamicButton.Name, parentFrame
        Case "pan_OSM":      mod_UI_Events.HandleOSMPanelEvents DynamicButton.Name, parentFrame
        Case "pan_Open":     mod_UI_Events.HandleOpenPanelEvents DynamicButton.Name, parentFrame
        ' . . . (pan_Closed, pan_PreviousMonth, pan_SelectPreviousMonth, ....)
    End Select
End Sub
\end{lstlisting}


% ==========================================
% BÖLÜM 4.3: MOD_CONFIG (GLOBAL YAPILANDIRMA VE SABİTLER)
% ==========================================

\section{\texttt{mod\_Config}}

\texttt{mod\_Config}, uygulamanın tüm merkezi sabitlerini, metinsel ifadelerini, renk paletini ve veritabanı sütun eşlemelerini bünyesinde barındıran standart kod modülüdür (\textit{Standard Module}). Kod karmaşıklığını önlemek ve bakımı kolaylaştırmak adına proje içerisindeki hiçbir fonksiyona sert kodlanmış (\textit{Hardcoded}) metin veya sayısal değer yazılmamıştır.

Gelecekte projede yapılabilecek güncellemeler doğrudan bu modül üzerinden yönetilir:

\begin{itemize}\setlength{\itemsep}{0pt}\setlength{\parskip}{1pt}
    \item \textbf{Arayüz Metinleri ve Açılır Pencereler (Pop-up):} Arayüzdeki buton yazıları (\texttt{BTN\_CAPTION}), panel başlıkları (\texttt{LBL\_TITLE}) veya kullanıcıya gösterilen tüm uyarı/bilgi bildirim mesajları (\texttt{MSG\_*}) değiştirilmek istendiğinde yalnızca bu modüldeki ilgili sabitler güncellenir.
    \item \textbf{Veritabanı ve Sütun Yapısı Değişiklikleri:} İlerleyen süreçte \texttt{ARIZA}, \texttt{BAKIM} veya \texttt{GECENAY} sayfalarında sütun yerleri veya aranan metin bağlamları değişirse; \texttt{DB\_*}, \texttt{CLOSED\_*} ve \texttt{PREV\_MONTH\_*} grubu altında tanımlanan sütun harfleri (\texttt{"A"}, \texttt{"D"}, \texttt{"K"} vb.) kolayca revize edilebilir.
\end{itemize}

\vspace{0.15cm}

\begin{infobox}{\faCogs\ Merkezi Yönetim ve Bakım Kolaylığı}
Uygulamanın renk paleti (\texttt{COLOR\_GREEN}, \texttt{COLOR\_RED}), panel boyutları (\texttt{PANEL\_WIDTH}) ve dinamik buton referanslarını tutan global \texttt{ButtonCollection} koleksiyonu bu modülde tanımlanarak projenin modülerliği garanti altına alınmıştır.
\end{infobox}

\vspace{0.15cm}

Aşağıdaki kod bloğu, \texttt{mod\_Config} modülünün kategorize edilmiş temsilî yapısını göstermektedir:

\begin{lstlisting}[caption={mod\_Config Modülü Sabit ve Yapılandırma Tanımları}]
'*****************************************************************************
' Module: mod_Config
' Purpose: Global configuration settings, constants, and variables.
'*****************************************************************************

' 1. SAYFA VE VERİTABANI İSİMLERİ (ARIZA, BAKIM, GECENAY)
Public Const MAIN_SHEET_NAME           As String = "Yenile"
Public Const PREVIOUS_MONTH_SHEET_NAME As String = "GECENAY"
Public Const DATABASE_SHEET_NAME_ARIZA As String = "ARIZA"
Public Const DATABASE_SHEET_NAME_BAKIM As String = "BAKIM"

' 2. VERİTABANI SÜTUN EŞLEMELERİ (COLUMN MAP)
Public Const DB_OSM_DATES_START_COL As String = "A"
Public Const DB_OPEN_CELL_COL       As String = "D"
Public Const DB_CLOSED_NAME_COL     As String = "B"
Public Const ALL_CASES_ASSIGNED_COL As String = "K"
' . . . (Diğer sütun ve hücre adres tanımlamaları)

' 3. POP-UP BİLDİRİM VE UYARI MESAJLARI (MESSAGES)
Public Const MSG_OPEN_SAVE_SUCCESS   As String = "Açık kayıt sayıları veritabanına kaydedildi."
Public Const MSG_ALL_OPERATIONS_COMPLETED As String = "Tüm güncelleme işlemleri başarıyla tamamlandı!"
' . . . (Diğer bildirim ve hata mesajları)

' 4. GLOBAL DEĞİŞKENLER VE ETKİLEŞİM KOLEKSİYONU
Public ButtonCollection As New Collection
\end{lstlisting}


\newpage

% ==========================================
% BÖLÜM 4.4: MOD_UI_MANAGER (PANEL VE ARAYÜZ YÖNETİMİ)
% ==========================================

\section{\texttt{mod\_UI\_Manager}}

\texttt{mod\_UI\_Manager}, dinamik panel geçişlerini, OLE nesnelerinin bellek yönetimini ve oturum kilit durumlarını orkestre eden servis katmanıdır. Sayfa üzerindeki tüm arayüz döngüsünün orkestrasyonunu üstlenir.

\vspace{0.3cm}

\textbf{\textcolor{primary}{1. \texttt{InitializeAppUI} Prosedürü}}

Çalışma kitabı açıldığında veya oturum sıfırlandığında tetiklenen ilk ilklendirme giriş noktasıdır. Uygulama altyapısını sırayla ayağa kaldırır.

\begin{lstlisting}[caption={InitializeAppUI Prosedürü Kod Yapısı}]
Public Sub InitializeAppUI()
    If Not Application.ActiveProtectedViewWindow Is Nothing Then Exit Sub
    Dim ws As Worksheet
    On Error Resume Next
    Set ws = ThisWorkbook.Sheets(MAIN_SHEET_NAME)
    On Error GoTo 0
    If ws Is Nothing Then Exit Sub
    
    BuildNavigationSystem
    SwitchTo "Login"
    
    With ws
        .ScrollArea = "A1"
        .EnableSelection = xlNoSelection
    End With
End Sub
\end{lstlisting}

\begin{infobox}{\faPlay\ Fonksiyon Özeti}
Navigasyon altyapısını kurup açılış ekranı olarak \texttt{Login} panelini çağırır. Kullanıcının sayfa hücreleriyle doğrudan etkileşime girmesini engellemek adına \texttt{ScrollArea} ve seçim izinlerini kısıtlar.
\end{infobox}

\vspace{0.3cm}

\textbf{\textcolor{primary}{2. \texttt{BuildNavigationSystem} Prosedürü}}

Tüm arayüz panellerini sıfırdan oluşturan ve bellek üzerindeki eski nesne referanslarını güvenle temizleyen yapıcı kısımdır.

\begin{lstlisting}[caption={BuildNavigationSystem Prosedürü (Özet)}]
Public Sub BuildNavigationSystem()
    ' Dynamic Button Koleksiyonunu ve OLE Nesnelerini Temizleme
    If Not ButtonCollection Is Nothing Then
        Dim btnItem As Object
        For Each btnItem In ButtonCollection
            Set btnItem.DynamicButton = Nothing
        Next btnItem
        Set ButtonCollection = Nothing
    End If
    Set ButtonCollection = New Collection
    
    ' . . . (bx_Anchor kontrolü ve eski OLE nesnelerinin temizlenmesi)
    
    ' Panellerin mod_UI_View Üzerinden Çizilmesi
    mod_UI_View.CreateLoginPanel ws
    mod_UI_View.CreateOSMPanel ws
    mod_UI_View.CreateOpenPanel ws
    ' . . . (Closed, PreviousMonth vb. diğer panellerin oluşturulması)
End Sub
\end{lstlisting}

\newpage

\textbf{\textcolor{primary}{3. \texttt{SwitchTo} Prosedürü}}

Paneller arası görünürlük geçişini ve liste bileşenlerinin dinamik veriye bağlanmasını sağlayan ana yönlendirme fonksiyonudur.

\begin{lstlisting}[caption={SwitchTo Prosedürü (Özet)}]
Public Sub SwitchTo(panelName As String)
    Dim ws As Worksheet, oleObj As OLEObject, realFrame As MSForms.Frame
    Set ws = ThisWorkbook.Sheets(MAIN_SHEET_NAME)
    If ws.OLEObjects("pan_" & panelName) Is Nothing Then Exit Sub
    
    ' Tüm Panelleri Gizle ve Hedef Paneli Konumlandır
    ' . . . (ws.OLEObjects("pan_*").Visible = False ve hedef panel hizalama)

    Set realFrame = oleObj.Object
    
    ' Panele Özel ListBox Formatlaması ve Render Yönlendirmesi
    Select Case panelName
        Case "OSM"
            ' . . . (lst_Status genişlik/sütun ayarları)
            mod_UI_View.RenderOSMSheetStatus realFrame, False
        Case "Open"
            ' . . . (lst_OpenStatus genişlik/sütun ayarları)
            mod_UI_View.RenderOpenSheets realFrame
        ' . . . (Closed, PreviousMonth vb. diğer Case yapıları)
    End Select
End Sub
\end{lstlisting}

\begin{infobox}{\faCog\ Fonksiyon Özeti}
Diğer servis ve event modüllerinden tek satırla çağrılarak yumuşak panel geçişi sağlar. Ekrandaki tüm panelleri gizleyip hedef paneli hizalar, \texttt{ListBox} sütun ayarlarını yapar ve ilgili \texttt{Render} prosedürünü tetikler.
\end{infobox}

\vspace{0.3cm}

\textbf{\textcolor{primary}{4. Yardımcı Kilit Fonksiyonları (\texttt{LockPanel} vb.)}}

İş akışı sırasında geçersiz durumlarda kullanıcının hatalı adım atmasını önlemek amacıyla buton erişimlerini kısıtlayan yardımcı prosedürlerdir.

\begin{infobox}{\faLock\ Fonksiyon Özeti}
\texttt{LockPanel} ve \texttt{LockPreviousMonthPanel} prosedürleri; işlem adımları tamamlandığında ilgili panellerdeki "Sonraki Adım" butonlarının \texttt{Enabled} özelliğini \texttt{False} konumuna getirip renklerini griye dönüştürerek geçişleri kilitler.
\end{infobox}

\vspace{0.4cm}


\begin{tcolorbox}[colback=yellow!15,colframe=orange!80!black,leftrule=3.5pt,boxrule=0.4pt,top=6pt,bottom=6pt]
    \textbf{\textcolor{orange!80!black}{\faInfoCircle\ mod\_UI\_Manager Genel Özet ve Mimari Rolü:}}\\
    \small \texttt{mod\_UI\_Manager}, arayüzün dinamik yaşam döngüsünü tamamen arka plandaki iş mantığından (\textit{Business Logic}) ve çizim bileşenlerinden (\textit{View Layer}) soyutlar. 
    
    Tüm panel oluşturma, temizleme, boyutlandırma ve erişim kısıtlama süreçlerini tek bir merkezde toplayarak projenin Modüler Yönetim, Bellek Güvenliği ve Kullanıcı Deneyimi (UX) standartlarını teminat altına alır.
\end{tcolorbox}


\newpage

% ==========================================
% BÖLÜM 4.5: MOD_UI_EVENTS (ARAYÜZ OLAY VE ETKİLEŞİM HANDLER KATMANI)
% ==========================================

\section{\texttt{mod\_UI\_Events}}

\texttt{mod\_UI\_Events}, kullanıcı arayüzündeki (UI) etkileşimleri, tıklama olaylarını (\textit{Events}) ve form girdilerini karşılayarak arka plandaki iş mantığına (\textit{Business Logic}) aktaran olay işleyici (\textit{Event Handler}) servis katmanıdır.

Sınıf modülleri (\texttt{clsDynamicButton}) üzerinden gelen tıklama sinyalleri, butonun içerisinde bulunduğu üst panel bağlamına (\texttt{MSForms.Frame}) göre bu modüldeki ilgili prosedürlere aktarılır. Modül, gelen parametreleri ve mevcut uygulama durumunu değerlendirerek yönlendirme ve veri işleme adımlarını yürütür.

Modülün üstlendiği temel kritik görevler şunlardır:

\begin{itemize}\setlength{\itemsep}{0pt}\setlength{\parskip}{0pt}
    \item \textbf{Dinamik Akış ve Yönlendirme Yönetimi:} Kullanıcının panel geçiş taleplerinde ("İleri", "Geri") mevcut iş durumunu kontrol ederek doğru panellere geçişi sağlar.
    \item \textbf{Veri Güncelleme Operasyonlarının Tetiklenmesi:} "Güncelle" butonlarına basıldığında, aktif panel üzerindeki veri bileşenlerini okur ve yazma operasyonu için \texttt{mod\_Services} katmanına iletir.
    \item \textbf{Kullanıcı Girdisi Doğrulama (Validation):} Geçmiş ay seçimi ekranlarında girilen verilerin tiplerini denetleyerek hatalı veri girişini engeller.
    \item \textbf{Arayüz Durum Güncellemeleri:} İşlem tamamlandığında ilgili butonların renklerini ve erişim izinlerini dinamik olarak günceller.
\end{itemize}

\vspace{0.15cm}

\begin{infobox}{\faHandPointer\ Mimari Soyutlama ve Trafik Kontrol Rolü}
\texttt{mod\_UI\_Events} modülü bünyesinde arayüz çizim kodlarını veya veritabanı okuma/yazma algoritmalarını barındırmaz. Sadece kullanıcı isteklerini doğru servislere ileten bir \textbf{Trafik Kontrolörü} gibi çalışır.
\end{infobox}

\vspace{0.15cm}

\textbf{\textcolor{primary}{1. Olay Karşılama ve Akış Denetimi Mekanizması}}

Aşağıdaki kod parçası, panel olaylarının (\texttt{HandleLoginPanelEvents}, \texttt{HandleSelectPreviousMonthEvents}) nasıl yakalandığını ve iş akışına göre yönlendirildiğini göstermektedir:

\begin{lstlisting}[caption={Panel Geçiş ve Akış Kontrol Mantığı}]
Public Sub HandleLoginPanelEvents(buttonName As String, parentFrame As MSForms.Frame)
    If buttonName = "btn_ToOSM" Then
        Dim hasOpen As Boolean, hasClosed As Boolean
        hasOpen = mod_Services.HasAnyOpenCase()
        hasClosed = mod_Services.HasAnyClosedCase()
        
        If Not hasOpen And Not hasClosed Then
            MsgBox MSG_NO_SHEETS_TO_PROCESS, vbInformation, MSG_TITLE_INFO
        ElseIf Not hasOpen And hasClosed Then
            mod_UI_Manager.SwitchTo "Closed"
        Else
            mod_UI_Manager.SwitchTo "OSM"
        End If
    ElseIf buttonName = "btn_ToPreviousMonth" Then
        ' ... Geckmis ay tarama kontrolleri ...
    End If
End Sub
\end{lstlisting}

\newpage

\begin{lstlisting}[caption={Girdi Doğrulama ve İkincil Akış Yönetimi}]
Public Sub HandleSelectPreviousMonthEvents(buttonName As String, parentFrame As MSForms.Frame)
    Dim txtYear As MSForms.TextBox, cmbMonth As MSForms.ComboBox
    
    If buttonName = "btn_NextFromSelectMonth" Then
        Set txtYear = parentFrame.Controls("txt_YearInput")
        Set cmbMonth = parentFrame.Controls("cmb_MonthSelect")
        
        If cmbMonth.ListIndex <> -1 And IsNumeric(txtYear.Text) And Len(Trim(txtYear.Text)) = 4 Then
            SelectedPreviousYear = CInt(txtYear.Text)
            SelectedPreviousMonth = cmbMonth.ListIndex + 1
            mod_UI_Manager.SwitchTo "PrevMonthClosed"
        Else
            MsgBox MSG_INVALID_YEAR_INPUT, vbExclamation, MSG_TITLE_ERROR
        End If
    End If
End Sub
\end{lstlisting}

\vspace{0.3cm}

\textbf{\textcolor{primary}{2. Veritabanı Yazma Tetikleme ve Görsel Güncelleme}}

Veri kaydetme butonlarına basıldığında çalışan özel fonksiyonlar (\texttt{HandleUpdateOpen}) ve işlem sonrasında buton görsellerini güncelleyen yardımcı prosedür aşağıdadır:

\begin{lstlisting}[caption={Veri Yazma Tetikleme ve Buton Durum Güncelleme}]
Private Sub HandleUpdateOpen()
    Dim ws As Worksheet, realFrame As MSForms.Frame
    Set ws = ThisWorkbook.Sheets(MAIN_SHEET_NAME)
    Set realFrame = ws.OLEObjects("pan_Open").Object
    
    If mod_Services.ExecuteDatabaseWrite(realFrame.Controls("lst_OpenStatus"), cellColIdx:=4, valColIdx:=3, sheetNameColIdx:=0) Then
        SetUpdateButtonSuccess realFrame.Controls("updateOpen"), realFrame.Controls("nextOpenButton")
        MsgBox MSG_OPEN_SAVE_SUCCESS, vbInformation, MSG_TITLE_SUCCESS
    End If
End Sub

Private Sub SetUpdateButtonSuccess(btnUpdate As MSForms.CommandButton, btnNext As MSForms.CommandButton)
    btnUpdate.BackColor = RGB(46, 204, 113)
    btnUpdate.Caption = BTN_CAPTION_UPDATED
    btnUpdate.Enabled = False
    btnNext.Enabled = True
    btnNext.BackColor = COLOR_GREEN
    btnNext.ForeColor = COLOR_WHITE
End Sub
\end{lstlisting}

\vspace{0.4cm}


\begin{tcolorbox}[colback=yellow!15,colframe=orange!80!black,leftrule=3.5pt,boxrule=0.4pt,top=5pt,bottom=5pt]
    \textbf{\textcolor{orange!80!black}{\faInfoCircle\ mod\_UI\_Events Genel Özet ve Mimari Rolü:}}\\
    \small \texttt{mod\_UI\_Events}, arayüz ile iş katmanı arasında gevşek bağlılık (\textit{Loose Coupling}) ilkesini uygular. Olay bazlı tetiklemeleri merkezi bir yapıda toplayarak uygulamanın bakımını kolaylaştırır, kod tekrarlarını önler ve kullanıcı etkileşimlerinin güvenli bir iş akışı çerçevesinde gerçekleşmesini teminat altına alır.
\end{tcolorbox}